\begin{itemframe}{مقدمه}
\itm
الگوریتم‌های ترتیبی
\fn{serial algorithms}
برای اجرا روی یک پردازنده طراحی شده‌اند.
این فصل، مدل الگوریتمی ما را گسترش می‌دهد تا الگوریتم‌های موازی
\fn{parallel algorithms}
را نیز در بر بگیرد، که در آن‌ها چندین دستور می‌توانند به طور هم‌زمان اجرا شوند.

\itm
 به طور خاص، ما مدل الگوریتم‌های وظیفه-موازی
\fn{task-parallel}
را بررسی می‌کنیم.
ساده‌ترین نوع الگوریتم‌های وظیفه-موازی، الگوریتم‌های انشعاب-الحاق
\fn{fork-join}
 نام دارد که ما در این فصل این نوع را مطالعه می‌کنیم.
\end{itemframe}


\begin{itemframe}{مقدمه}
\itm
الگوریتم‌های انشعاب-الحاق را می‌توان با استفاده از افزودن امکانات ساده‌ایی به همان کد ترتیبی عادی، به‌صورت واضح بیان کرد. این الگوریتم‌ها در عمل نیز می‌توانند به‌صورت کارآمد پیاده‌سازی شوند.
\itm
به کامپیوترهایی که بیش‌ از یک پردازنده داشته باشند چندهسته‌ایی
\fn{multicore}
به هر پردازنده یک هسته گفته می‌شود و هر هسته می‌تواند به حافظه مشترک
\fn{shared memory}
کامپیوتر دسترسی داشته باشد.
\itm
یکی از روش‌های برنامه‌نویسی کامپیوترهای چند‌هسته‌ای، موازی‌سازی نخ‌محور
\fn{thread parallelism}
است.
این مدل از یک انتزاع نرم‌افزاری به نام «پردازنده‌ی مجازی» یا نخ
\fn{thread}
استفاده می‌کند. نخ‌ها همگی حافظه‌ی مشترکی را به اشتراک می‌گذارند. هر نخ شمارنده‌ی
\fn{program counter}
اختصاصی خود را دارد و می‌تواند به‌طور مستقل از نخ‌های دیگر اجرا شود.
\end{itemframe}


\begin{itemframe-s}{مقدمه}{برنامه‌نویسی وظیفه-موازی}
\itm
اما برنامه‌نویسی نخ‌محور دشوار و مستعد خطاست. برای مثال تقسیم‌کار میان نخ‌ها به گونه‌ای که هر نخ بار کاری تقریباً برابری دریافت کند، ممکن است پیچیده باشد. بنابراین به سراغ مدل وظیفه-موازی می‌رویم.
\itm
در این مدل یک لایه‌ی نرم‌افزاری روی نخ‌ها قرار می‌دهیم تا پردازنده‌های یک کامپیوتر چند‌هسته‌ای را مدیریت کند. به این لایه نرم افزاری سکوی وظیفه-موازی
\fn{task-parallel platform}
گفته می‌شود که می‌تواند یک کتابخانه یا یک زبان برنامه‌نویسی باشد.
\itm
در برنامه‌نویسی وظیفه-موازی، برنامه‌نویس مشخص می‌کند که چه وظایف محاسباتی‌ای ممکن است به‌صورت موازی اجرا شوند، اما مشخص نمی‌کند که کدام نخ یا پردازنده، آن وظیفه را اجرا کند.
\end{itemframe-s}


\begin{itemframe-s}{مقدمه}{برنامه‌نویسی وظیفه-موازی}
\itm
سکوی وظیفه-موازی دارای یک زمان‌بند است که به طور متوازن وظایف را میان پردازنده‌ها تقسیم می‌کند.
\itm
پس برنامه‌نویس درگیر پروتکل‌های ارتباطی، توازن بار کاری، و سایر پیچیدگی‌های برنامه‌نویسی با نخ نمی‌شود.
\itm
الگوریتم‌هایی وظیفه موازی
\fn{Task-parallel algorithms}
 نیز تعمیم یافته همان الگوریتم‌های سریالی هستند.
\end{itemframe-s}


\begin{itemframe-s}{مقدمه}{موازی‌سازی انشعاب-‌الحاق}
\itm
تقریباً همه‌ی محیط‌های وظیفه-موازی، از موازی‌سازی انشعاب-الحاق پشتیبانی می‌کنند که معمولاً در دو ویژگی زبانی ظاهر می‌شود: انشعاب
\fn{spawning}
و حلقه موازی
\fn{parallel loop} .
\itm
انشعاب یعنی بعد از فراخوانی یک تابع، فراخوانی‌کننده به اجرای خود ادامه دهند. در حالی که  تابع فراخوانی شده نیز مشغول محاسبه‌ی نتیجه‌اش است.
\itm
حلقه‌ی موازی مانند یک حلقه‌ی for معمولی است، با این تفاوت که چندین تکرار از حلقه می‌توانند به‌طور هم‌زمان اجرا شوند.
\end{itemframe-s}


\begin{itemframe-s}{مقدمه}{موازی‌سازی انشعاب-‌الحاق}
\itm
برخلاف مدل چند نخی برنامه نویس در اینجا مشخص نمی‌کند که کدام وظایف «باید» حتما به‌صورت موازی اجرا شوند، بلکه تنها مشخص می‌کند که کدام وظایف «می‌توانند» به‌صورت موازی اجرا شوند. این ویژگی از مدل وظیفه-موازی به ارث برده شده.

\itm
موازی‌سازی انشعاب-الحاق چندین مزیت دارد:
\item[1]
این مدل برنامه‌نویسی تعمیم یافته مدل برنامه‌نویسی ترتیبی است. که با آن آشنایی کامل داریم.
برای توصیف یک الگوریتم موازی انشعاب-هم‌گرا، فقط سه واژه‌ی کلیدی به شبه‌کدها اضافه می‌شود:
$$parallel, spawn, sync$$
\end{itemframe-s}


\begin{itemframe-s}{مقدمه}{موازی‌سازی انشعاب-‌الحاق}
\itm
حذف این واژه‌های کلیدی از شبه‌کد موازی منجر به ساخت شبه‌کد ترتیبی برای همان مسئله می‌شود که آن را «طرح ترتیبی»
\fn{serial projection}
 الگوریتم موازی می‌نامیم.
\item[2]
مدل وظیفه-موازی راهی دقیق برای سنجش میزان موازی‌سازی فراهم می‌آورد.
\item[3]
با استفاده از قابلیت انشعاب بسیاری از الگوریتم‌های تقسیم‌و‌حل به سادگی تبدیل به الگوریتم‌های موازی می‌شوند. همچنین، همانند الگوریتم‌های ترتیبی تقسیم‌و‌حل، الگوریتم‌های موازی در مدل انشعاب-الحاق نیز توسط رابطه‌های بازگشتی قابل تحلیل اند.
\item[4]
انشعاب-الحاق تنها یک مدل تئوری نیست و تقریباً همه‌ٔ محیط‌های چند‌هسته‌ای از آن پشتیبانی می‌کنند.
\end{itemframe-s}